%&pdflatex
\documentclass[a4paper,12pt]{article}

\addtolength{\textwidth}{1cm}

\usepackage[spanish]{babel}

\usepackage[leqno]{amsmath}
\usepackage{amsfonts,amsthm}

\nofiles %% .aux files unnecessary

\usepackage{setspace}
\newcommand{\tab}{\hspace{5mm}}

\theoremstyle{plain}
\newtheorem*{theoremOne}{Teorema 1}
\newtheorem*{propTwo}{Proposici\'on 2}
\newtheorem*{corThree}{Corolario 3}
\newtheorem*{corFour}{Corolario 4}
\newtheorem*{theoremFive}{Teorema 5}

\theoremstyle{definition}
\newtheorem*{note}{Nota}

\newcommand{\N}{\mathbb{N}}

\date{}

%%%%%%%%%%%%%%%%%%%%%%%%%%%%%%%%%%%%%%%%%%%%%%%%%%%%%%%%%%%%%%%
\title{\large \bfseries \boldmath ACOTACI\'ON EN $L^{p}$ DE LOS N\'UCLEOS
RELATIVOS A PESOS DE JACOBI GENERALIZADOS}

\author{%
J. J. Guadalupe Hern\'andez$^{(1)}$, M. P\'{e}rez Riera$^{(1)}$,\\[2pt]
F. J. Ruiz Blasco$^{(1)}$ y J. L. Varona Malumbres$^{(2)}$\\[5pt]
\small $^{(1)}$ Departamento de Matem\'aticas,\\[-2pt]
\small Facultad de Ciencias, Universidad de Zaragoza.\\[3pt]
\small $^{(2)}$ Departamento de Matem\'atica Aplicada,\\[-2pt]
\small Colegio Universitario de La Rioja, Universidad de Zaragoza.%
\renewcommand\thefootnote{}%
\renewcommand\footnotemark{}%
\footnote{Este art\'{\i}culo es una traducci\'on al espa\~nol de
\textsc{J.~J.~Guadalupe, M.~P\'{e}rez,
F.~J.~Ruiz y J.~L.~Varona},
{$L^{p}$-boundedness of the kernels relative
to generalized Jacobi weights}, en
``Polinomios Ortogonales y Aplicaciones''
(Actas del VI Simposium, Gij\'on, 1989),
168--178, Universidad de Oviedo, Gij\'on, 1990.}%
}
%%%%%%%%%%%%%%%%%%%%%%%%%%%%%%%%%%%%%%%%%%%%%%%%%%%%%%%%%%%%%%%

\begin{document}

%%%%%%%%%%%%%%%%%%%%%%%%%%%%%%%%%%%%%%%%%%%%
\maketitle
%%%%%%%%%%%%%%%%%%%%%%%%%%%%%%%%%%%%%%%%%%%%

%

\noindent\textit{Clasificaci\'on A.M.S.}: 42C10.

\begin{abstract}
Sea $w$ un peso de Jacobi generalizado sobre el intervalo $[-1,1]$,
es decir, $w(x) = h(x)(1-x)^{\alpha} (1+x)^{\beta}
\prod_{i=1}^{N} |x-x_{i}|^{\gamma_{i}}$, con $\alpha, \beta,
\gamma_{i} > -1$ y ciertas condiciones de continuidad
sobre~$h$. Mediante el uso de la teor\'{\i}a de pesos $A_{p}$, se puede
demostrar la convergencia de la serie de Fourier de polinomios
ortonormales con respecto a $w$ cuando $\gamma_{i} \geq 0$
$\forall i$. En este trabajo obtenemos acotaciones
de las normas en $L^{p}$ de los n\'ucleos relativos a $w$, que permiten
extender el resultado de la convergencia en media al caso general.
\end{abstract}

%%%%%%%%%%%%%%%%%%%%%%%%%%%%%%%%%%%%%%%%%%%%

Sea $w$ un peso de Jacobi generalizado, es decir,
$$
w(x) = h(x)(1-x)^{\alpha} (1+x)^{\beta}
\prod_{i=1}^{N} |x-x_{i}|^{\gamma_{i}},
\qquad x \in [-1,1],
$$
donde:
\begin{itemize}
\item[(a)] $\alpha, \beta, \gamma_{i} > -1$,
$t_{i} \in (-1,1)$, $t_{i}\neq t_{j}$ $\forall i \neq j$;
\item[(b)] $h$ es una funci\'on continua y positiva en $[-1,1]$ y
$\omega(h,\delta) \delta^{-1} \in L^{1}(0,2)$,
siendo $\omega(h,\delta)$ el m\'odulo de continuidad de~$h$.
\end{itemize}

Sea $d\mu = w(x)\,dx$ sobre $[-1,1]$ y $S_{n}$ ($n \geq 0$)
la suma parcial en\'{e}sima de la serie de Fourier de los polinomios
ortonormales con respecto a $d\mu$. El estudio de la acotaci\'on
\begin{equation}
\| S_{n}f \|_{L^{p}(u^{p}\,d\mu)}
\leq C \| f \|_{L^{p}(v^{p}\,d\mu)},
\tag{1}
\end{equation}
donde
$$
u(x) = (1- x)^{a}(1+ x)^{b} \prod_{i=1}^{N} | x- t_{i} |^{g_{i}}
$$
y
$$
v(x) = (1- x)^{A}(1+ x)^{B} \prod_{i=1}^{N} | x- t_{i} |^{G_{i}}
$$
fue realizado por Badkov ([B]) para el caso $u = v$
mediante la obtenci\'on directa de estimaciones de los n\'ucleos
$K_{n}(x,y)$ asociados a los polinomios ortonormales con respecto a~$d\mu$.
Posteriormente, Varona ([V]) consider\'o el mismo
problema, con los pesos $u$ y $v$ no no necesariamente iguales; utilizando
otro m\'{e}todo, que consiste en hacer uso de la teor\'{\i}a de pesos
$A_{p}$, encontr\'o condiciones para la acotaci\'on (1) que
generalizan las obtenidas para $u = v$ por Badkov.
Sin embargo, el resultado, que enunciamos a continuaci\'on,
se deduce s\'olo en el caso $\gamma_{i} \geq 0$ $\forall i$:

\begin{theoremOne}
Sea $\gamma_{i} \geq 0$ $\forall i$ y $1 < p < \infty$.
Si se cumplen las desigualdades
\begin{equation}
\left\{
\begin{aligned}
A + (\alpha+1) \left(\frac{1}{p} - \frac{1}{2}\right) &< \frac{1}{4}, \\
B + (\beta+1) \left(\frac{1}{p} - \frac{1}{2}\right) &< \frac{1}{4}, \\
G_{i} + (\gamma_{i}+1) \left(\frac{1}{p} - \frac{1}{2}\right) &< \frac{1}{2}
\quad \forall i,
\end{aligned}
\right.
\tag{2}
\end{equation}
\begin{equation}
\left\{
\begin{aligned}
A + (\alpha+1) \left(\frac{1}{p} - \frac{1}{2}\right) &< \frac{\alpha+1}{2}, \\
B + (\beta+1) \left(\frac{1}{p} - \frac{1}{2}\right) &< \frac{\beta+1}{2}, \\
G_{i} + (\gamma_{i}+1) \left(\frac{1}{p} - \frac{1}{2}\right) &< \frac{\gamma_{i}+1}{2}
\quad \forall i,
\end{aligned}
\right.
\tag{3}
\end{equation}
\begin{equation}
\left\{
\begin{aligned}
a + (\alpha+1) \left(\frac{1}{p} - \frac{1}{2}\right) &> - \frac{1}{4}, \\
b + (\beta+1) \left(\frac{1}{p} - \frac{1}{2}\right) &> - \frac{1}{4}, \\
g_{i} + (\gamma_{i}+1) \left(\frac{1}{p} - \frac{1}{2}\right) &> - \frac{1}{2}
\quad \forall i,
\end{aligned}
\right.
\tag{4}
\end{equation}
\begin{equation}
\left\{
\begin{aligned}
a + (\alpha+1) \left(\frac{1}{p} - \frac{1}{2}\right) &> - \frac{\alpha+1}{2}, \\
b + (\beta+1) \left(\frac{1}{p} - \frac{1}{2}\right) &> - \frac{\beta+1}{2}, \\
g_{i} + (\gamma_{i}+1) \left(\frac{1}{p} - \frac{1}{2}\right) &> - \frac{\gamma_{i}+1}{2}
\quad \forall i,
\end{aligned}
\right.
\tag{5}
\end{equation}
\begin{equation}
A \leq a, \qquad B \leq b, \qquad G_{i} \leq g_{i}
\quad \forall i,
\tag{6}
\end{equation}
entonces existe $C > 0$ tal que
$$
\| S_{n}f \|_{L^{p}(u^{p}\,d\mu)}
\leq C \| f \|_{L^{p}(v^{p}\,d\mu)}
\quad \forall f \in L^{p}(v^{p}\,d\mu),
\quad \forall n \in \N.
$$
\end{theoremOne}

El objeto de este trabajo es eliminar la restricci\'on $\gamma_{i} \geq 0$,
mediante el estudio de las normas de los n\'ucleos $K_{n}(x,y)$.
Utilizaremos la siguiente notaci\'on: $\{P_{n}(x)\}$, $\{k_{n}\}$ y
$\{K_{n}(x,y)\}$ ser\'an, respectivamente, los polinomios
ortonormales, sus coeficientes directores y los n\'ucleos relativos
a~$d\mu$; si $c \in (-1,1)$, $\{P_{n}^{c}(x)\}$,
$\{k_{n}^{c}\}$ y $\{K_{n}^{c}(x,y)\}$ ser\'an
los correspondientes a $(x-c)^{2}\,d\mu$.
No es dif\'{\i}cil entonces establecer las siguientes relaciones,
$\forall n \geq 1$:
\begin{equation}
K_{n}(x,y) = (x-c) (y-c) K_{n-1}^{c}(x,y)
+ \frac{K_{n}(x,c)K_{n}(c,y)}{K_{n}(c,c)}\,,
\tag{7}
\end{equation}
\begin{equation}
K_{n}(x,c) = \frac{k_{n}}{k_{n}^{c}}\, P_{n}(c) P_{n}^{c}(x)
- \frac{k_{n-1}^{c}}{k_{n+1}}\, P_{n+1}(c) P_{n-1}^{c}(x).
\tag{8}
\end{equation}
Puede demostrarse adem\'as (v\'{e}ase [MNT], teoremas~10 y~11,
y~[R], p\'ag.~212) que
\begin{equation}
\lim_{n \to \infty} \frac{k_{n}}{k_{n}^{c}}
= \lim_{n \to \infty} \frac{k_{n-1}^{c}}{k_{n+1}}
= \frac{1}{2}.
\tag{9}
\end{equation}

Si definimos
$$
d(x,n) = (1 - x + n^{-2})^{-(2\alpha+1)/4}
(1 + x + n^{-2})^{-(2\beta+1)/4}
\prod_{i=1}^{N} ( | x-t_{i}| + n^{-1})^{-\gamma_{i}/2},
$$
es conocido ([B]) que existe una constante $C$ tal que
$\forall x \in [-1,1]$, $\forall n$,
\begin{equation}
| P_{n}(x) | \leq C d(x,n).
\tag{10}
\end{equation}

Tambi\'{e}n se conocen estimaciones de los n\'ucleos, una de las
cuales es \'esta ([N], p\'ag.~4 y p\'ag.~119, teorema~25): si
$c \in (-1,1)$ y el factor $| x - c |$
est\'a presente en $w$ con un exponente $\gamma$, existen constantes
positivas $C_{1}$ y $C_{2}$, que dependen de $c$,
tales que,~$\forall n$,
\begin{equation}
C_{1} n^{\gamma+1} \leq K_{n}(c,c) \leq C_{2} n^{\gamma+1}.
\tag{11}
\end{equation}

A partir de ahora designaremos por $C$ todas las constantes, con
lo que $C$ ser\'a una constante posiblemente distinta cada vez
que aparezca. Utilizando (8), (9) y (10), llegamos al siguiente
resultado:

\begin{propTwo}
Sea $1 < p < \infty$, $1/p + 1/q = 1$.
Supongamos que se cumplen las desigualdades (4) y (5). Sea $-1 < c < 1$
y sean $\gamma$ y $g$ los exponentes de $| x - c |$
en $w$ y $u$, respectivamente. Entonces existe una constante positiva
$C$ tal que, para todo $n \geq 0$, se tiene:
$$
\| K_{n}(x,c) \|_{L^{p}(u^{p}w)}
\leq
\begin{cases}
C n^{(\gamma+1)/q - g}, &\text{si}\quad g < (\gamma +1) (1/2 - 1/p) + 1/2,\\
C n^{\gamma/2} (\log n)^{1/p}, &\text{si}\quad g = (\gamma +1) (1/2 - 1/p) + 1/2,\\
C n^{\gamma/2}, &\text{si}\quad (\gamma +1) (1/2 - 1/p) + 1/2 < g.
\end{cases}
$$
\end{propTwo}

\begin{proof}
De (10) se deduce que $| P_{n}(c) | \leq C n^{\gamma/2}$.
Como $\{P_{n}^{c}\}$ es la sucesi\'on
asociada a $(x-c)^{2}\,d\mu$, tambi\'{e}n de (10) se deduce que
$$
| P_{n}^{c}(x) | \leq C ( | x - c | + n^{-1})^{-1} d(x,n).
$$
A partir de (8) y (9) obtenemos, entonces,
\begin{equation}
| K_{n}(x,c) | \leq C n^{\gamma/2} ( | x - c | + n^{-1})^{-1} d(x,n).
\tag{12}
\end{equation}
Fijemos ahora $\varepsilon > 0$ tal que $| t_{i} - c | > \varepsilon$
$\forall i$ que haga $t_{i} \neq c$.
Podemos poner
\begin{multline*}
\| K_{n}(x,c) \|^{p}_{L^{p}(u^{p}w)} \\
= \int_{| x - c | > \varepsilon} | K_{n}(x,c) |^{p} u(x)^{p} w(x)\,dx
+ \int_{| x - c | < \varepsilon} | K_{n}(x,c) |^{p} u(x)^{p} w(x)\,dx.
\end{multline*}
Utilizando la cota (12) obtenemos, para el primer sumando:
\begin{align*}
&\int_{| x - c | > \varepsilon} | K_{n}(x,c) |^{p} u(x)^{p} w(x)\,dx \\
&\qquad\leq C n^{p \gamma/2} \int_{| x - c | > \varepsilon}
( | x - c | + n^{-1})^{-p} d(x,n)^{p} u(x)^{p} w(x)\,dx \\
&\qquad\leq C_{1} n^{p \gamma/2} \int_{-1}^{1} d(x,n)^{p} u(x)^{p} w(x)\,dx.
\end{align*}
De (4) y (5) es f\'acil deducir que esta \'ultima integral est\'a
acotada por una constante que no depende de $n$, con lo que
\begin{equation}
\int_{| x - c | > \varepsilon} | K_{n}(x,c) |^{p} u(x)^{p} w(x)\,dx
\leq C n^{p\gamma/2}.
\tag{13}
\end{equation}
Veamos ahora el segundo sumando; como en $| x - c | < \varepsilon$
existe una constante $C$ tal que, $\forall n$,
$d(x,n) \leq C ( | x - c | + n^{-1})^{-\gamma/2}$,
$u(x) \leq C | x - c |^{g}$
y $w(x) \leq C | x - c |^{\gamma}$, tenemos
\begin{align*}
&\int_{| x - c | < \varepsilon} | K_{n}(x,c) |^{p} u(x)^{p} w(x)\,dx \\
&\qquad\leq C n^{p\gamma/2} \int_{| x - c | < \varepsilon} ( | x - c | + n^{-1})^{-p}
d(x,n)^{p} u(x)^{p} w(x)\,dx \\
&\qquad\leq C_{1} n^{p \gamma/2}
\int_{| x - c | < \varepsilon} ( | x - c | + n^{-1})^{-p(1+\gamma/2)}
| x - c |^{gp+\gamma} \,dx \\
&\qquad\leq C _{2} n^{p\gamma/2}
\int_{0}^{1} (y + n^{-1})^{-p(1+\gamma/2)} y^{gp+\gamma} \,dy \\
&\qquad= C_{2} n^{p\gamma/2+p(1+\gamma/2)-gp-\gamma-1}
\int_{0}^{1} (ny + 1)^{-p(1+\gamma/2)} (ny)^{gp+\gamma} n\,dy \\
&\qquad= C_{2} n^{p\gamma/2+p(1+\gamma/2)-gp-\gamma-1}
\int_{0}^{n} (r + 1)^{-p(1+\gamma/2)} r^{gp+\gamma} \,dr.
\end{align*}
Observando que
$p (1+\gamma/2) - gp - \gamma - 1 = p [ (\gamma+1) (1/2 - 1/p) - g + 1/2]$
y que existen constantes $C_{1}$ y $C_{2}$ tales que
$C_{1} \leq r+1 \leq C_{2}$ en $[0,1]$ y $C_{1}r \leq r+1 \leq C_{2}r$
en $[1,n]$, llegamos finalmente a la desigualdad
\begin{equation}
\begin{aligned}
&\int_{| x - c | < \varepsilon} | K_{n}(x,c) |^{p} u(x)^{p} w(x)\,dx \\
&\qquad\leq C n^{p \gamma/2+p [ (\gamma+1)(1/2 - 1/p) - g+1/2 ]}
\int_{0}^{1} r^{gp+\gamma} \,dr \\
&\qquad\qquad\null+ C n^{p\gamma/2+p [ (\gamma+1)(1/2 - 1/p) - g+1/2 ]}
\int_{1}^{n} r^{-p [ (\gamma+1)(1/2 - 1/p) - g+1/2 ] -1} \,dr.
\end{aligned}
\tag{14}
\end{equation}
Como (5) implica $gp + \gamma > -1$,
acotamos el primer sumando de esta forma:
\begin{equation}
C n^{p\gamma/2+p [ (\gamma+1)(1/2 - 1/p) - g+1/2 ]}
\int_{0}^{1} r^{gp+\gamma} \,dr
\leq C n^{p \gamma/2+p [ (\gamma+1)(1/2 - 1/p) - g+1/2 ]}.
\tag{15}
\end{equation}
Distingamos ahora los tres casos del enunciado:

    (a) Si $g < (\gamma +1) (1/2 - 1/p) + 1/2$,
entonces $-p [ (\gamma +1) (1/2 - 1/p) -g+ 1/2 ] -1 < -1$, con lo que
$$
\int_{1}^{n} r^{-p [ (\gamma+1)(1/2 - 1/p) - g+1/2 ] -1} \,dr \leq C.
$$
De (14) y (15) se sigue, en este caso,
$$
\int_{| x - c | < \varepsilon} | K_{n}(x,c) |^{p} u(x)^{p} w(x)\,dx
\leq C n^{p\gamma/2+p [ (\gamma+1)(1/2 - 1/p) - g+1/2 ]}.
$$
Puesto que $p [ (\gamma +1)(1/2 - 1/p) - g+1/2 ] > 0$,
de esta desigualdad y (13) se obtiene
\begin{align*}
\| K_{n}(x,c) \|^{p}_{L^{p}(u^{p}w)}
&\leq C n^{p\gamma/2+p [ (\gamma+1)(1/2 - 1/p) - g+1/2 ]} \\
&= C n^{p [ (\gamma+1)(1 - 1/p) - g ]}
= C n^{p [ (\gamma+1)/q - g ]},
\end{align*}
como quer\'{\i}amos probar.

    (b) Si $(\gamma +1) (1/2 - 1/p) + 1/2 < g$,
entonces $-p [ (\gamma +1) (1/2 - 1/p) -g+ 1/2 ] -1 > -1$, con lo que
$$
\int_{1}^{n} r^{-p [ (\gamma+1)(1/2 - 1/p) - g+1/2 ] -1} \,dr
\leq C n^{-p [ (\gamma+1)(1/2 - 1/p) - g+1/2 ] }.
$$
De (14) y (15) se sigue, ahora,
$\int_{| x - c | < \varepsilon} | K_{n}(x,c) |^{p} u(x)^{p} w(x)\,dx
\leq C n^{p\gamma/2}$
y, finalmente,
$$
\| K_{n}(x,c) \|^{p}_{L^{p}(u^{p}w)} \leq C n^{p\gamma/2}.
$$

    (c) Por \'ultimo, si $g = (\gamma +1) (1/2 - 1/p) + 1/2$,
$$
\int_{1}^{n} r^{-p [ (\gamma+1)(1/2 - 1/p) - g+1/2 ] -1} \,dr \leq C \log n;
$$
de donde,
$$
\int_{| x - c | < \varepsilon} | K_{n}(x,c) |^{p} u(x)^{p} w(x)\,dx
\leq C n^{p\gamma/2} \log n
$$
y
$$
\| K_{n}(x,c) \|^{p}_{L^{p}(u^{p}w)} \leq C n^{p\gamma/2} \log n.
$$
De esta manera, queda demostrada la proposici\'on.
\end{proof}

\begin{corThree}
Sea $1 < p < \infty$, $1/p + 1/q = 1$.
Supongamos que se cumplen las desigualdades (2) y (3). Sea $-1 < c < 1$
y sean $\gamma$ y $G$ los exponentes de $| x - c |$
en $w$ y $v$, respectivamente. Entonces existe una constante positiva
$C$ tal que, para todo $n \geq 0$, se tiene:
$$
\| K_{n}(x,c) \|_{L^{q}(v^{-q}w)}
\leq
\begin{cases}
C n^{\gamma/2}, &\text{si}\quad G < (\gamma +1) (1/2 - 1/p) - 1/2,\\
C n^{\gamma/2} (\log n)^{1/q}, &\text{si}\quad G = (\gamma +1) (1/2 - 1/p) - 1/2,\\
C n^{(\gamma+1)/p + G}, &\text{si}\quad (\gamma +1) (1/2 - 1/p) - 1/2 < G.
\end{cases}
$$
\end{corThree}

\begin{proof}
Basta aplicar la proposici\'on~2 al peso $v^{-1}$ y tener en cuenta
que $1/2 - 1/p = 1/q - 1/2$.
\end{proof}

El siguiente resultado es el que necesitamos para ampliar el
teorema~1 al caso general $\gamma_{i} > -1$:

\begin{corFour}
Sea $1 < p < \infty$, $1/p + 1/q = 1$.
Supongamos que se cumplen las desigualdades (2), (3), (4), (5)
y (6). Sea $-1 < c < 1$.
Entonces existe una constante positiva $C$ tal que, para todo $n \geq 0$,
se tiene:
$$
\| K_{n}(x,c) \|_{L^{p}(u^{p}w)} \| K_{n}(x,c) \|_{L^{q}(v^{-q}w)}
\leq C K_{n}(c,c).
$$
\end{corFour}

\begin{proof}
Basta aplicar, caso por caso, la proposici\'on~2, el corolario~3
y la estimaci\'on~(11).
\end{proof}

\begin{note}
Aaunque en lo que sigue no lo vamos a utilizar, el corolario~4 es cierto
tambi\'{e}n en el caso $c = \pm 1$.
La demostraci\'on es an\'aloga: partiendo de otras expresiones
para $K_{n}(x,\pm 1)$, pueden probarse
resultados similares a la proposici\'on~2 y el corolario~3, de
donde se deduce el corolario~4.
\end{note}

Ahora estamos en condiciones de generalizar el teorema~1 al caso
$\gamma_{i} > -1$, como hab\'{\i}amos dicho:

\begin{theoremFive}
Sea $1 < p < \infty$.
Si se verifican las desigualdades
$$
\left\{
\begin{aligned}
A + (\alpha+1) \left(\frac{1}{p} - \frac{1}{2}\right) &< \frac{1}{4}, \\
B + (\beta+1) \left(\frac{1}{p} - \frac{1}{2}\right) &< \frac{1}{4}, \\
G_{i} + (\gamma_{i}+1) \left(\frac{1}{p} - \frac{1}{2}\right) &< \frac{1}{2}
\quad \forall i,
\end{aligned}
\right.
\quad
\left\{
\begin{aligned}
A + (\alpha+1) \left(\frac{1}{p} - \frac{1}{2}\right) &< \frac{\alpha+1}{2}, \\
B + (\beta+1) \left(\frac{1}{p} - \frac{1}{2}\right) &< \frac{\beta+1}{2}, \\
G_{i} + (\gamma_{i}+1) \left(\frac{1}{p} - \frac{1}{2}\right) &< \frac{\gamma_{i}+1}{2}
\quad \forall i,
\end{aligned}
\right.
$$
$$
\left\{
\begin{aligned}
a + (\alpha+1) \left(\frac{1}{p} - \frac{1}{2}\right) &> - \frac{1}{4}, \\
b + (\beta+1) \left(\frac{1}{p} - \frac{1}{2}\right) &> - \frac{1}{4}, \\
g_{i} + (\gamma_{i}+1) \left(\frac{1}{p} - \frac{1}{2}\right) &> - \frac{1}{2}
\quad \forall i,
\end{aligned}
\right.
\quad
\left\{
\begin{aligned}
a + (\alpha+1) \left(\frac{1}{p} - \frac{1}{2}\right) &> - \frac{\alpha+1}{2}, \\
b + (\beta+1) \left(\frac{1}{p} - \frac{1}{2}\right) &> - \frac{\beta+1}{2}, \\
g_{i} + (\gamma_{i}+1) \left(\frac{1}{p} - \frac{1}{2}\right) &> - \frac{\gamma_{i}+1}{2}
\quad \forall i,
\end{aligned}
\right.
$$
$$
A \leq a, \qquad B \leq b, \qquad G_{i} \leq g_{i}
\quad \forall i,
$$
entonces existe $C > 0$ tal que
$$
\| S_{n}f \|_{L^{p}(u^{p}\,d\mu)}
\leq C \| f \|_{L^{p}(v^{p}\,d\mu)}
\quad \forall f \in L^{p}(v^{p}\,d\mu),
\quad \forall n \in \N.
$$
\end{theoremFive}

\begin{proof}
Por inducci\'on sobre el n\'umero de exponentes $\gamma_{i}$ negativos.
Si $\gamma_{i} \geq 0$ $\forall i$, el resultado es cierto,
como hemos comentado anteriormente (teorema~1).

Supongamos que hay $k$ exponentes $\gamma_{i}$ negativos,
con $k > 0$, y que el resultado es cierto para $k-1$. Sea $c \in (-1,1)$
un punto con exponente $\gamma$ negativo.
Recordemos la f\'ormula~(7):
$$
K_{n}(x,y) = (x-c) (y-c) K_{n-1}^{c}(x,y)
+ \frac{K_{n}(x,c) K_{n}(c,y)}{K_{n}(c,c)}.
$$
Definamos los operadores
\begin{align*}
T_{n}f(x) &= \int_{-1}^{1}
\frac{K_{n}(x,c) K_{n}(c,y)}{K_{n}(c,c)}\, f(y) w(y)\,dy,\\
R_{n}f(x) &= \int_{-1}^{1}
(x-c) (y-c) K_{n-1}^{c}(x,y) f(y) w(y)\,dy.
\end{align*}
Entonces, $S_{n} = T_{n} + R_{n}$. Estudiamos primero
los operadores~$T_{n}$:
$$
T_{n}f(x) = \int{K_{n}(x,c)}{K_{n}(c,c)}
\int_{-1}^{1} K_{n}(c,y) f(y) w(y)\,dy,
$$
luego
\begin{align*}
\| T_{n}f \|_{L^{p}(u^{p}w)}
&\leq \frac{\int_{-1}^{1}  | K_{n}(c,y) | v(y)^{-1} | f(y) | v(y) w(y) \,dy}
{K_{n}(c,c)}\, \| K_{n}(x,c) \|_{L^{p}(u^{p}w)} \\
&\leq \frac{\| K_{n}(x,c) \|_{L^{p}(u^{p}w)}
\| K_{n}(x,c) v(x)^{-1} \|_{L^{q}(w)}}{K_{n}(c,c)}\,
\| f v \|_{L^{p}(w)} \\
&= \frac{\| K_{n}(x,c) \|_{L^{p}(u^{p}w)}
\| K_{n}(x,c) \|_{L^{q}(v^{-q}w)}}{K_{n}(c,c)}\,
\| f \|_{L^{p}(v^{p}w)}.
\end{align*}
Del corolario~4 se deduce entonces que $\| T_{n}f \|_{L^{p}(u^{p}\,d\mu)}
\leq C \| f \|_{L^{p}(v^{p}\,d\mu)}$
$\forall f \in L^{p}(v^{p}\,d\mu)$,
$\forall n \in \N$. S\'olo hace falta probar ahora esta misma
acotaci\'on para los operadores~$R_{n}$.

Ahora bien, si por $S_{n}^{c}$ denotamos las sumas parciales
de la serie de Fourier con respecto a la medida $(x-c)^{2}w(x)\,dx$,
resulta
$$
R_{n}f(x) = (x-c) \int_{-1}^{1} (y-c) K_{n-1}^{c}(x,y) f(y) w(y)\,dy
= (x-c) S_{n-1}^{c} \left(\frac{f(y)}{y-c}, x \right),
$$
luego
\begin{align*}
&\| R_{n}f \|_{L^{p}(u^{p}w)} \leq C \| f \|_{L^{p}(v^{p}w)}
\quad\forall f \in L^{p}(v^{p}w) \\
&\qquad\iff
\big\| (x-c) S_{n-1}^{c} \big(\tfrac{f(y)}{y-c}, x \big) \big\|_{L^{p}(u^{p}w)}
\leq C \| f \|_{L^{p}(v^{p}w)}
\quad\forall f \in L^{p}(v^{p}w) \\
&\qquad\iff
\| (x-c) S_{n-1}^{c}g (x) \|_{L^{p}(u^{p}w)}
\leq C \| (x-c) g \|_{L^{p}(v^{p}w)}
\quad\forall g \in L^{p}(|x-c|^{p} v^{p}w) \\
&\qquad\iff
\| S_{n-1}^{c}g \|_{L^{p}(|x-c|^{p}u^{p}w)}
\leq C \| g \|_{L^{p}(|x-c|^{p}v^{p}w)}
\quad\forall g \in L^{p}(|x-c|^{p} v^{p}w) \\
&\qquad\iff
\| S_{n-1}^{c}g \|_{L^{p}({\tilde{u}}^{p}(x-c)^{2}w)}
\leq C \| g \|_{L^{p}({\tilde{v}}^{p}(x-c)^{2}w)}
\quad\forall g \in L^{p}({\tilde{v}}^{p}(x-c)^{2}w),
\end{align*}
donde $\tilde{u}(x) = | x - c |^{1-2/p} u(x)$
y $\tilde{v}(x) = | x - c |^{1-2/p} v(x)$.

Se trata por lo tanto de la acotaci\'on de las sumas parciales
$S_{n}^{c}$ con el par de pesos $( \tilde{u}, \tilde{v} )$.
Pero la serie de Fourier que estamos considerando ahora se refiere
al peso de Jacobi generalizado $(x-c)^{2}w(x)$, que tiene s\'olo $k-1$
exponentes $\gamma_{i}$ negativos, puesto que en el
punto $c$ el exponente es $\gamma + 2 > 1$.
Por hip\'otesis, el teorema es cierto en este caso y solamente
tenemos que comprobar que se verifican las condiciones del enunciado,
pero para los pesos
$(x-c)^{2}w(x)$, $| x - c |^{1-2/p} u(x)$
y $| x - c |^{1-2/p} v(x)$.

Salvo en el punto $c$, estos pesos tienen los mismos exponentes
que $w$, $u$ y $v$, luego dichas condiciones son las mismas y, por
lo tanto se verifican. En el punto $c$, los exponentes son, respectivamente,
$$
\gamma + 2,\qquad g + 1 - 2/p,\qquad G + 1 - 2/p.
$$
As\'{\i} que las desigualdades que debemos comprobar son:
\begin{align*}
(G + 1 - 2/p) + ( \gamma + 2 + 1 ) (1/p - 1/2 ) &< 1/2,\\
(G + 1 - 2/p) + ( \gamma + 2 + 1 ) (1/p - 1/2 ) &< ( \gamma + 2 + 1 )/2,\\
(g + 1 - 2/p) + ( \gamma + 2 + 1 ) (1/p - 1/2 ) &> - 1/2,\\
(g + 1 - 2/p) + ( \gamma + 2 + 1 ) (1/p - 1/2 ) &> - ( \gamma + 2 + 1 )/2,\\
G + 1 - 2/p &\leq g + 1 - 2/p.
\end{align*}
Es inmediato ver que todas ellas se verifican, a partir de
las hip\'otesis del enunciado. Por consiguiente, se cumple la
acotaci\'on
$$
\| S_{n-1}^{c}g \|_{L^{p}({\tilde{u}}^{p}(x-c)^{2}w)}
\leq C \| g \|_{L^{p}({\tilde{v}}^{p}(x-c)^{2}w)}
\quad\forall g \in L^{p}({\tilde{v}}^{p}(x-c)^{2}w),
\quad\forall n \in \N;
$$
luego tambi\'{e}n
$$
\| R_{n}f \|_{L^{p}(u^{p}w)} \leq C \| f \|_{L^{p}(v^{p}w)}
\quad\forall f \in L^{p}(v^{p}w), \quad\forall n \in \N.
$$
Y, por lo tanto,
$$
\| S_{n}f \|_{L^{p}(u^{p}\,d\mu)} \leq C \| f \|_{L^{p}(v^{p}\,d\mu)}
\quad\forall f \in L^{p}(v^{p}\,d\mu), \quad\forall n \in \N.
$$
Es decir, el resultado es v\'alido para $k$ exponentes $\gamma_{i}$ negativos.
Por inducci\'on, es v\'alido en general y el teorema
queda demostrado.
\end{proof}

\begin{note}
Puede demostrarse que el rec\'{\i}proco tambi\'{e}n es
cierto, es decir, que la acotaci\'on de las sumas parciales $S_{n}$
implica las cinco condiciones del teorema.
\end{note}


%%%%%%%%%%%%%%%%%%%%%%%%%%%%%%%%%%
\begin{thebibliography}{MNT}
%%%%%%%%%%%%%%%%%%%%%%%%%%%%%%%%%%

\frenchspacing

\bibitem[B]{B}
\textsc{V. M. Badkov},
Convergence in the mean and almost everywhere of Fourier series
in polynomials orthogonal on an interval,
\textit{Math USSR Sb.} \textbf{24} (1974), 223--256.

\bibitem[MNT]{MNT}
\textsc{A. M\'at\'{e}, P. Nevai, V. Totik},
Extensions of Szeg\H{o}'s theory of orthogonal polynomials.~II,
\textit{Constr. Approx.} \textbf{3} (1987), 51--72.

\bibitem[N]{N}
\textsc{P. Nevai},
``Orthogonal Polynomials'',
Memoirs of the Amer. Math. Soc., vol.~18,
n\textsuperscript{o}~213, Providence, R.I., U.S.A., 1979.

\bibitem[R]{R}
\textsc{E. A. Rahmanov},
On the asymptotics of the ratio of orthogonal polynomials,
\textit{Math. USSR Sb.} \textbf{32} (1977), 199--213.

\bibitem[V]{V}
\textsc{J. L. Varona},
``Convergencia en L$^{p}$ con pesos de la serie de Fourier
respecto de algunos sistemas ortogonales'', Tesis doctoral,
Sem. Mat. Garc\'{\i}a de Galdeano, sec.~2,
n.\textsuperscript{o}~22. Zaragoza, 1989.

\end{thebibliography}

\end{document}
